
\documentclass[11pt]{article}

\usepackage[a4paper,margin=2.5cm]{geometry}
\usepackage[utf8]{inputenc}
\usepackage[T1]{fontenc}
\usepackage{amsmath,amssymb,amsfonts}
\usepackage{graphicx}
\usepackage{booktabs}
\usepackage{hyperref}
\usepackage{authblk}
\usepackage{bm}
\usepackage{microtype}
\usepackage{xcolor}

\hypersetup{
  colorlinks=true,
  linkcolor=blue!60!black,
  citecolor=blue!60!black,
  urlcolor=blue!60!black
}

\newcommand{\DC}{D_C}
\newcommand{\Geff}{G_{\rm eff}}
\newcommand{\sig}{\sigma_8}
\newcommand{\LCDM}{\Lambda\mathrm{CDM}}
\newcommand{\BuP}{\textsc{BuP}}
\newcommand{\camb}{\textsc{camb}}
\newcommand{\ent}{\rm ent}

\title{
  \textbf{\'Equation d'Einstein effective \'emergente\\
  dans la Gravit\'e Quantique Bottom-Up (\BuP)}
}

\author{Farid Hamdad}

\affil{
  Cadre Bottom-Up Quantum Gravity (\BuP)\\
  \small\url{https://github.com/Farid-Hamdad/Bottom-up-Quantum-Gravity}
}

\date{2026}

\begin{document}

\maketitle

% ─────────────────────────────────────────────────────────────
\begin{abstract}
Dans \BuP, la relativit\'e g\'en\'erale n'est pas postul\'ee comme une dynamique fondamentale
de l'espace-temps ; elle appara\^{\i}t comme la limite continue et thermodynamique de la
r\'eponse du r\'eseau d'intrication aux variations locales d'information mutuelle.

Nous proposons une \'equation d'Einstein effective \'emergente,
\[
G_{\mu\nu}[g^{\ent}]
=
8\pi\Geff(d_s)\,T_{\mu\nu}^{\rm matter}
+
8\pi G\,T_{\mu\nu}^{\ent}[d_s],
\]
o\`u $g^{\ent}_{\mu\nu}$ est la m\'etrique \'emergente issue de la distance d'intrication,
$d_s$ la dimension spectrale locale,
\[
\Geff(d_s)=\frac{2}{d_s-1}G,
\]
et $T_{\mu\nu}^{\ent}[d_s]$ un tenseur effectif construit \`a partir des gradients et
variations de la dimension spectrale.

Aucune constante cosmologique $\Lambda$ n'est postul\'ee. Aucune composante d'\'energie
sombre n'est ajout\'ee. L'acc\'el\'eration tardive de l'Univers est interpr\'et\'ee comme
une cons\'equence g\'eom\'etrique de l'\'evolution de la dimension spatiale effective
$d_{\rm bg}(z) < 3$ dans le r\'egime cosmologique homog\`ene.

Cette \'equation g\'en\'eralise les r\'esultats cosmologiques ant\'erieurs du programme
\BuP\ : elle reproduit la dynamique homog\`ene $d_{\rm bg}(z)$, la r\'esolution de la
tension $\sigma_8$, et la r\'eduction dimensionnelle locale dans les surdensit\'es
proto-galactiques. Sa limite $d_s\rightarrow 3$, $\nabla_\mu d_s\rightarrow 0$ redonne
la relativit\'e g\'en\'erale standard.
\end{abstract}

% ─────────────────────────────────────────────────────────────
\section{Introduction}
% ─────────────────────────────────────────────────────────────

Les travaux pr\'ec\'edents du programme Bottom-Up Quantum Gravity (\BuP) ont \'etabli
qu'une g\'eom\'etrie effective peut \^etre reconstruite \`a partir de la structure
d'intrication d'un \'etat quantique global~\cite{hamdad2026bup}. Dans ce cadre, la
distance, la dimension effective et le couplage gravitationnel ne sont pas postul\'es
comme des donn\'ees fondamentales : ils \'emergent du r\'eseau d'information mutuelle
entre sous-syst\`emes quantiques.

Trois r\'egimes compl\'ementaires ont d\'ej\`a \'et\'e explor\'es.

\textbf{Fond cosmologique \`a dimension variable.}
Une dimension homog\`ene $d_{\rm bg}(z)$, variable au cours du temps cosmologique,
permet de d\'ecrire l'\'evolution tardive de la g\'eom\'etrie effective. Ce mod\`ele,
contraint par les donn\'ees BAO et supernovæ, am\'eliore l'ajustement par rapport
\`a $\LCDM$ tout en pr\'eservant la physique du CMB par sa valeur critique asymptotique
\[
d_{\rm bg}(z\to\infty)=\DC\simeq 3.0598.
\]

\textbf{Croissance lin\'eaire et tension $\sigma_8$.}
Le couplage gravitationnel effectif
\[
\Geff(z)=\frac{2}{d_{\rm bg}(z)-1}G,
\]
appliqu\'e de fa\c{c}on coh\'erente \`a l'expansion du fond et aux perturbations
lin\'eaires, conduit \`a $\sig\simeq 0.772$, en accord avec les contraintes de
lentillage faible, sans param\`etre suppl\'ementaire~\cite{hamdad2026p3}.

\textbf{R\'egime non lin\'eaire et observations JWST.}
L'extension locale
\[
d_{\rm local}(z,\delta)=d_{\rm bg}(z)-\varepsilon\delta
\]
d\'ecrit une r\'eduction de la dimension spectrale dans les r\'egions surdenses.
Ce m\'ecanisme, dont l'existence microscopique est v\'erifi\'ee sur des graphes
d'intrication finis (\(\varepsilon_{\rm micro}(N=20)=0.0065\)),
acc\'el\`ere l'effondrement des halos massifs \`a haut redshift et fournit une
explication semi-quantitative de l'exc\`es de galaxies massives observ\'e par
JWST~\cite{hamdad2026p4}.

Ces r\'esultats sugg\`erent une image unifi\'ee : la gravitation macroscopique
serait la r\'eponse g\'eom\'etrique effective du r\'eseau d'intrication aux
variations locales de connectivit\'e informationnelle.

Cependant, une \'etape conceptuelle essentielle manquait encore : formuler
l'\'equation dynamique qui relie la courbure \'emergente aux sources mat\'erielles
et informationnelles. Dans la relativit\'e g\'en\'erale, cette relation est donn\'ee
par les \'equations d'Einstein,
\[
G_{\mu\nu}+\Lambda g_{\mu\nu}=8\pi G\,T_{\mu\nu}.
\]
Dans \BuP, l'objectif n'est pas de postuler cette \'equation comme fondamentale,
mais de montrer qu'une forme effective de type Einstein appara\^{\i}t dans la
limite continue du graphe d'intrication.

\textbf{Objectif de ce travail.}
Nous proposons une \'equation d'Einstein effective \'emergente,
\[
G_{\mu\nu}[g^{\ent}]
=
8\pi\Geff(d_s)\,T_{\mu\nu}^{\rm matter}
+
8\pi G\,T_{\mu\nu}^{\ent}[d_s],
\]
o\`u $g^{\ent}_{\mu\nu}$ est la m\'etrique reconstruite depuis la distance
d'intrication, $d_s(x)$ la dimension spectrale locale, et
$\Geff(d_s)=\frac{2}{d_s-1}G$.

Dans \BuP, la relativit\'e g\'en\'erale n'est pas postul\'ee comme une dynamique
fondamentale de l'espace-temps ; elle appara\^{\i}t comme la limite continue et
thermodynamique de la r\'eponse du r\'eseau d'intrication aux variations locales
d'information mutuelle.

Cette \'equation g\'en\'eralise les r\'esultats ant\'erieurs :
\begin{itemize}
  \item dans la limite homog\`ene, elle redonne la cosmologie \`a dimension variable ;
  \item dans le r\'egime lin\'eaire, elle reproduit la r\'esolution de la tension $\sig$ ;
  \item dans les r\'egions surdenses, elle retrouve le m\'ecanisme local de r\'eduction
        dimensionnelle.
\end{itemize}
Sa limite $d_s\to 3$, $\nabla_\mu d_s\to 0$ redonne la relativit\'e g\'en\'erale standard.

\textbf{Structure du papier.}
La Section~\ref{sec:metric} rappelle la construction de la m\'etrique \'emergente
$g_{\mu\nu}^{\ent}$ \`a partir de la matrice d'information mutuelle.
La Section~\ref{sec:curvature} introduit la courbure effective et la dimension
spectrale locale $d_s(x)$.
La Section~\ref{sec:closure} justifie la fermeture effective depuis le r\'eseau
d'intrication.
La Section~\ref{sec:einstein} pr\'esente l'\'equation d'Einstein effective \BuP,
son tenseur d'intrication, sa conservation covariante et ses limites physiques.
La Section~\ref{sec:discussion} discute les perspectives observationnelles et
num\'eriques.
La Section~\ref{sec:conclusion} conclut.

% ─────────────────────────────────────────────────────────────
\section{M\'etrique \'emergente d'intrication}
\label{sec:metric}
% ─────────────────────────────────────────────────────────────

Dans \BuP, la distance spatiale n'est pas une donn\'ee primitive. Elle est
reconstruite \`a partir de la structure d'intrication d'un \'etat quantique global.
Le point de d\'epart est la matrice d'information mutuelle
\[
W_{ij}=I(i:j),
\]
o\`u $I(i:j)$ mesure la corr\'elation informationnelle entre deux sous-syst\`emes
$i$ et $j$.

L'id\'ee physique centrale est simple : deux degr\'es de libert\'e fortement intriqu\'es
doivent \^etre proches dans la g\'eom\'etrie \'emergente, tandis que deux degr\'es de
libert\'e faiblement corr\'el\'es doivent \^etre \'eloign\'es. On introduit donc une
longueur effective d'ar\^ete
\[
\ell_{ij}
=
\ell_0
\left[
-\log\!\left(
\frac{I(i:j)+\epsilon_I}{I_{\max}+\epsilon_I}
\right)
\right],
\]
o\`u $\ell_0$ est une \'echelle fondamentale de longueur, par exemple une longueur
de Planck effective. La quantit\'e $I_{\max}$ d\'esigne une \'echelle de normalisation
de l'information mutuelle, que l'on peut prendre comme le maximum
$\max_{i,j}I(i:j)$, ou plus prudemment comme une moyenne robuste sur les paires
effectivement connect\'ees du graphe. Le param\`etre $\epsilon_I$ est un r\'egulateur
positif.

Le r\^ole de $\epsilon_I$ est uniquement technique : il \'evite les divergences lorsque
$I(i:j)\rightarrow0$. On choisit $\epsilon_I \ll I(i:j)$ pour toutes les paires
effectivement connect\'ees du graphe. Ainsi, le choix pr\'ecis de $\epsilon_I$
n'affecte pas les distances entre paires non triviales ; il ne r\'egularise que les
liens d'information n\'egligeable.

Cette d\'efinition poss\`ede trois propri\'et\'es importantes :
\[
I(i:j)\simeq I_{\max}
\quad\Longrightarrow\quad
\ell_{ij}\simeq 0,
\]
\[
I(i:j)\ll I_{\max}
\quad\Longrightarrow\quad
\ell_{ij}\gg \ell_0,
\]
et
\[
I(i:j)\rightarrow 0
\quad\Longrightarrow\quad
\ell_{ij}\rightarrow \infty
\]
dans la limite $\epsilon_I\rightarrow0$. Ainsi, la g\'eom\'etrie \'emergente encode
directement la hi\'erarchie des corr\'elations quantiques.

La distance entre deux n\oe{}uds du r\'eseau d'intrication est d\'efinie comme la
distance g\'eod\'esique discr\`ete
\[
d_{ij}^{\rm ent}
=
\min_{\gamma:i\rightarrow j}
\sum_{(a,b)\in\gamma}
\ell_{ab},
\]
o\`u $\gamma$ parcourt tous les chemins reliant $i$ \`a $j$ dans le graphe pond\'er\'e.
Cette d\'efinition transforme la matrice d'information mutuelle en une structure
m\'etrique discr\`ete.

Dans la limite o\`u le nombre de degr\'es de libert\'e devient grand et o\`u la densit\'e
effective de n\oe{}uds reste suffisamment r\'eguli\`ere, on suppose que le graphe
d'intrication admet une limite continue au sens des r\'eseaux al\'eatoires
g\'eom\'etriques~\cite{penrose2003}. La distance discr\`ete converge alors vers la
distance g\'eod\'esique d'une m\'etrique effective $g_{\mu\nu}^{\rm ent}$ :
\[
d_{ij}^{\rm ent}
\longrightarrow
d_{g_{\rm ent}}(x_i,x_j),
\]
avec
\[
d_{g_{\rm ent}}(x_i,x_j)
=
\inf_{\gamma}
\int_{\gamma}
\sqrt{
g_{\mu\nu}^{\rm ent}(x)\,
dx^\mu dx^\nu
}.
\]

Les coordonn\'ees locales $x_i^\mu$ ne sont pas suppos\'ees donn\'ees a priori. Elles
peuvent \^etre reconstruites \`a partir de la matrice des distances $d_{ij}^{\rm ent}$
par des m\'ethodes de plongement spectral, de multidimensional scaling, ou de
plongement isom\'etrique local. La m\'etrique $g_{\mu\nu}^{\rm ent}$ est donc d\'efinie
op\'erationnellement comme la m\'etrique continue dont les distances g\'eod\'esiques
reproduisent, \`a grande \'echelle, les distances du graphe d'intrication.

Dans un voisinage suffisamment petit autour d'un point $x$, on approxime la
distance \'emergente par
\[
d_{ij}^{2}
\simeq
g_{\mu\nu}^{\rm ent}(x)
\Delta x_{ij}^{\mu}
\Delta x_{ij}^{\nu},
\]
o\`u $\Delta x_{ij}^{\mu}=x_i^\mu-x_j^\mu$. La m\'etrique locale est alors obtenue
par r\'egression locale des distances g\'eod\'esiques du graphe.

Cette construction justifie l'interpr\'etation de $g_{\mu\nu}^{\rm ent}$ comme
m\'etrique \'emergente. Les fluctuations locales de l'information mutuelle induisent
des fluctuations de longueur,
\[
\delta I(i:j)
\quad\Longrightarrow\quad
\delta \ell_{ij},
\]
puis des fluctuations de la m\'etrique effective,
\[
\delta \ell_{ij}
\quad\Longrightarrow\quad
\delta g_{\mu\nu}^{\rm ent}.
\]

Au premier ordre, une variation relative de l'information mutuelle donne
\[
\delta \ell_{ij}
=
-\ell_0
\frac{\delta I(i:j)}{I(i:j)+\epsilon_I}.
\]
Ainsi, une augmentation locale d'intrication raccourcit les distances \'emergentes,
tandis qu'une diminution locale d'intrication les allonge. Cette relation constitue
le m\'ecanisme g\'eom\'etrique \'el\'ementaire par lequel les variations d'intrication
produisent une courbure effective.

Dans la th\'eorie effective continue, ces variations microscopiques ne sont pas
d\'ecrites directement par $\delta I(i:j)$, mais par des observables g\'eom\'etriques
d\'eriv\'ees, en particulier la dimension spectrale locale $d_s(x)$. Autrement dit,
la variation d'intrication est m\'edi\'ee, dans la description effective, par les
variations spatiales et temporelles de $d_s(x)$.

La m\'etrique \'emergente peut \'egalement \^etre reli\'ee au Laplacien d'intrication
\[
L=D-W,
\qquad
D_{ii}=\sum_j W_{ij}.
\]
Lorsque la densit\'e de n\oe{}uds tend vers l'infini et que la limite continue existe,
le Laplacien discret converge vers le Laplace-Beltrami associ\'e \`a $g_{\mu\nu}^{\rm ent}$ :
\[
L
\;\longrightarrow\;
-\Delta_{g_{\rm ent}}
=
-\frac{1}{\sqrt{|g_{\rm ent}|}}
\partial_\mu
\!\left(
\sqrt{|g_{\rm ent}|}
g_{\rm ent}^{\mu\nu}
\partial_\nu
\right).
\]
Cette identification fournit un second crit\`ere de reconstruction : la m\'etrique
effective doit reproduire non seulement les distances g\'eod\'esiques, mais aussi les
propri\'et\'es spectrales du Laplacien d'intrication.

La dimension spectrale locale est alors d\'efinie par le noyau de chaleur
\[
P_x(\tau)
=
\left(e^{-\tau L}\right)_{xx},
\]
et
\[
d_s(x)
=
-2\frac{d\log P_x(\tau)}{d\log\tau}.
\]
Dans la limite continue, cette quantit\'e mesure la dimension effective explor\'ee
par la diffusion sur la g\'eom\'etrie $g_{\mu\nu}^{\rm ent}$.

La coh\'erence entre la distance \'emergente, le Laplacien d'intrication et la dimension
spectrale impose donc le sch\'ema suivant :
\[
I(i:j)
\quad\Longrightarrow\quad
\ell_{ij}
\quad\Longrightarrow\quad
d_{ij}^{\rm ent}
\quad\Longrightarrow\quad
g_{\mu\nu}^{\rm ent}
\quad\Longrightarrow\quad
d_s(x).
\]

Cette cha\^{\i}ne est le c\oe{}ur g\'eom\'etrique de \BuP. Elle montre que la m\'etrique
effective n'est pas une variable ind\'ependante : elle est une repr\'esentation continue
de la structure d'intrication sous-jacente. Par cons\'equent, le champ $d_s(x)$ n'est
pas une mati\`ere exotique ajout\'ee \`a la th\'eorie, mais la manifestation continue
des propri\'et\'es spectrales du r\'eseau d'intrication. Toute dynamique de l'information
mutuelle induit donc une dynamique de la m\'etrique \'emergente et de la dimension
spectrale.

% ─────────────────────────────────────────────────────────────
\section{Courbure effective et dimension spectrale}
\label{sec:curvature}
% ─────────────────────────────────────────────────────────────

La section pr\'ec\'edente d\'efinit la m\'etrique \'emergente $g_{\mu\nu}^{\rm ent}$
\`a partir des distances reconstruites sur le graphe d'intrication. Il reste maintenant
\`a d\'efinir la notion de courbure associ\'ee \`a cette g\'eom\'etrie. Dans \BuP, la
courbure n'est pas introduite comme une propri\'et\'e fondamentale de l'espace-temps :
elle appara\^{\i}t comme une mesure de l'inhomog\'en\'eit\'e de la connectivit\'e
informationnelle.

Le point de d\'epart est de nouveau le Laplacien d'intrication
\[
L = D-W,
\qquad
D_{ii}=\sum_j W_{ij},
\qquad
W_{ij}=I(i:j).
\]
Ce Laplacien encode \`a la fois la connectivit\'e locale, les propri\'et\'es de diffusion
et la structure spectrale du r\'eseau. Dans la limite continue, il converge vers
le Laplace-Beltrami associ\'e \`a la m\'etrique \'emergente :
\[
L\longrightarrow -\Delta_{g_{\rm ent}}.
\]
La g\'eom\'etrie effective peut donc \^etre sond\'ee par le noyau de chaleur
\[
K(x,x';\tau)
=
\left(e^{-\tau L}\right)_{x x'}.
\]

La trace diagonale du noyau de chaleur
\[
P_x(\tau)=K(x,x;\tau)
\]
mesure le retour local d'une diffusion fictive au point $x$ apr\`es un temps de
diffusion $\tau$. Elle d\'efinit la dimension spectrale locale :
\[
d_s(x,\tau)
=
-2\frac{d\log P_x(\tau)}{d\log\tau}.
\]
Lorsque la pente est approximativement constante sur une fen\^etre de diffusion, on
d\'efinit
\[
d_s(x)
\simeq
\left\langle
d_s(x,\tau)
\right\rangle_{\tau\in[\tau_{\min},\tau_{\max}]}.
\]
Dans les simulations \BuP\ actuellement accessibles, une fen\^etre typique
$\tau\in[0.01,50]$ donne une pente logarithmique suffisamment stable pour extraire
une dimension spectrale effective.

Cette d\'efinition est importante car elle ne suppose pas l'existence pr\'ealable
d'une dimension enti\`ere. Elle mesure directement la dimension effective explor\'ee
par la diffusion sur le graphe d'intrication. Une r\'egion fortement connect\'ee peut
donc pr\'esenter une dimension spectrale diff\'erente de la dimension topologique
apparente.

Dans une vari\'et\'e riemannienne continue de dimension $d$, le d\'eveloppement
asymptotique du noyau de chaleur est
\[
K(x,x;\tau)
\simeq
\frac{1}{(4\pi\tau)^{d/2}}
\left[
1+\frac{\tau}{6}R(x)+\mathcal{O}(\tau^2)
\right],
\]
o\`u $R(x)$ est le scalaire de courbure. Cette relation montre que les corrections
\`a la diffusion locale encodent directement la courbure.

Par analogie, dans \BuP, les \'ecarts du noyau de chaleur du graphe d'intrication
par rapport \`a une diffusion homog\`ene d\'efinissent une courbure effective. On peut
introduire un scalaire de courbure spectral
\[
R_{\rm ent}(x)
\sim
6\,
\frac{1}{\tau}
\left[
(4\pi\tau)^{d_s(x)/2}P_x(\tau)-1
\right],
\]
dans une fen\^etre de diffusion o\`u $d_s(x)$ varie lentement. Cette expression est
inspir\'ee du d\'eveloppement asymptotique du noyau de chaleur sur une vari\'et\'e,
o\`u elle est exacte \`a l'ordre lin\'eaire en $\tau$. Dans \BuP, elle sert de
d\'efinition op\'erationnelle de $R_{\rm ent}(x)$ sur une fen\^etre de diffusion o\`u
$d_s(x)$ est stable. Elle ne doit donc pas \^etre interpr\'et\'ee comme une \'egalit\'e
microscopique exacte, mais comme une prescription effective pour extraire une
courbure moyenne \`a partir du noyau de chaleur.

Une autre mesure naturelle de la courbure discr\`ete est la courbure d'Ollivier-Ricci
du graphe pond\'er\'e~\cite{ollivier2009}. Pour une ar\^ete $(i,j)$, elle est d\'efinie par
\[
\kappa_{ij}
=
1
-
\frac{W_1(m_i,m_j)}{d_{ij}^{\rm ent}},
\]
o\`u $W_1(m_i,m_j)$ est la distance de Wasserstein entre les mesures de voisinage
$m_i$ et $m_j$, et $d_{ij}^{\rm ent}$ la distance \'emergente entre les deux n\oe{}uds.
La mesure $m_i$ peut \^etre choisie comme la distribution de masse pond\'er\'ee sur le
voisinage de $i$,
\[
m_i(j)
=
\frac{W_{ij}}{\sum_k W_{ik}},
\]
ou, dans une approximation plus simple, comme une mesure uniforme sur les voisins
effectifs de $i$.

Dans la limite continue, la courbure d'Ollivier-Ricci le long d'une direction $v^\mu$
converge vers une contraction du tenseur de Ricci :
\[
\kappa_{ij}
\longrightarrow
R_{\mu\nu}^{\rm ent}v^\mu v^\nu.
\]
Cette relation fournit le pont entre la courbure discr\`ete du graphe d'intrication
et la courbure tensorielle de la m\'etrique \'emergente.

On dispose donc de deux diagnostics compl\'ementaires :
\[
P_x(\tau)
\quad\Longrightarrow\quad
d_s(x),\,R_{\rm ent}(x),
\]
et
\[
\kappa_{ij}
\quad\Longrightarrow\quad
R_{\mu\nu}^{\rm ent}.
\]
Le premier mesure la courbure moyenne via la diffusion locale ; le second mesure
la courbure directionnelle via le transport optimal sur le graphe. Leur coh\'erence
dans la limite continue constitue une condition forte pour l'\'emergence d'une
g\'eom\'etrie riemannienne effective.

Dans \BuP, la dimension spectrale $d_s(x)$ joue alors un r\^ole particulier. Elle
n'est pas seulement un diagnostic de dimension : elle devient une observable
g\'eom\'etrique locale contr\^olant la force gravitationnelle effective. Les r\'esultats
cosmologiques pr\'ec\'edents sugg\`erent la relation
\[
\Geff(d_s)
=
\frac{2}{d_s-1}G.
\]
Ainsi, une variation locale de $d_s$ modifie directement la r\'eponse gravitationnelle :
\[
\delta d_s(x)
\quad\Longrightarrow\quad
\delta\Geff(x).
\]

Au premier ordre,
\[
\frac{\delta\Geff}{\Geff}
=
-\frac{\delta d_s}{d_s-1}.
\]
Une diminution locale de la dimension spectrale augmente donc le couplage gravitationnel
effectif, tandis qu'une augmentation de $d_s$ le diminue. Ce m\'ecanisme relie
directement les fluctuations spectrales du graphe d'intrication \`a la dynamique
gravitationnelle \'emergente.

La cha\^{\i}ne g\'eom\'etrique compl\`ete devient alors
\[
\begin{aligned}
I(i:j)
&\quad\Longrightarrow\quad
L
\quad\Longrightarrow\quad
P_x(\tau)
\quad\Longrightarrow\quad
d_s(x)
\\
&\quad\Longrightarrow\quad
\Geff(d_s)
\quad\Longrightarrow\quad
\text{r\'eponse gravitationnelle}.
\end{aligned}
\]

Cette relation est le pont essentiel entre la g\'eom\'etrie d'intrication microscopique
et l'\'equation d'Einstein effective propos\'ee dans la section suivante.

% ─────────────────────────────────────────────────────────────
\section{Justification effective depuis le r\'eseau d'intrication}
\label{sec:closure}
% ─────────────────────────────────────────────────────────────

Nous montrons maintenant comment l'\'equation effective propos\'ee dans la section
suivante peut \^etre justifi\'ee \`a partir de la r\'eponse continue du r\'eseau
d'intrication. Cette justification n'a pas la pr\'etention d'\^etre une d\'erivation
microscopique compl\`ete ; elle \'etablit une fermeture effective minimale, coh\'erente
avec les propri\'et\'es g\'eom\'etriques et spectrales du r\'eseau.

\subsection{R\'eponse de la m\'etrique aux variations d'intrication}

Le point de d\'epart est le graphe pond\'er\'e $W_{ij}=I(i:j)$.
Une perturbation locale de l'\'etat quantique modifie les corr\'elations,
\[
I(i:j)\to I(i:j)+\delta I(i:j),
\]
et induit donc une variation des longueurs :
\[
\delta\ell_{ij}=-\ell_0\frac{\delta I(i:j)}{I(i:j)+\epsilon_I}+\mathcal{O}(\delta I^2).
\]
Dans la limite continue, les longueurs $\ell_{ij}$ d\'efinissent une m\'etrique effective
$g_{\mu\nu}^{\rm ent}$. Ainsi, une variation locale d'intrication produit une variation
de m\'etrique :
\[
\delta I(i:j)
\;\Longrightarrow\;
\delta\ell_{ij}
\;\Longrightarrow\;
\delta g_{\mu\nu}^{\rm ent}.
\]
La r\'eponse de la courbure \'emergente est alors donn\'ee, au premier ordre, par
\[
\delta G_{\mu\nu}[g^{\rm ent}]=D_{\mu\nu}^{\alpha\beta}\,\delta g_{\alpha\beta}^{\rm ent},
\]
o\`u $D_{\mu\nu}^{\alpha\beta}$ d\'esigne l'op\'erateur lin\'earis\'e d'Einstein autour
de la g\'eom\'etrie effective de fond.

\subsection{R\'eponse de la dimension spectrale}

La m\^eme perturbation de l'intrication modifie le Laplacien $L=D-W$, puis le
noyau de chaleur $P_x(\tau)=\left(e^{-\tau L}\right)_{xx}$, et donc la dimension
spectrale locale :
\[
\delta I(i:j)
\;\Longrightarrow\;
\delta L
\;\Longrightarrow\;
\delta P_x(\tau)
\;\Longrightarrow\;
\delta d_s(x).
\]
La dimension spectrale est donc l'observable continue qui encode la r\'eponse spectrale
du r\'eseau \`a une perturbation informationnelle. Les r\'esultats cosmologiques \BuP\
indiquent que la r\'eponse gravitationnelle locale est contr\^ol\'ee par
$\Geff(d_s)=\frac{2}{d_s-1}G$, d'o\`u au premier ordre
$\delta\Geff/\Geff=-\delta d_s/(d_s-1)$.

\subsection{Fermeture effective minimale}

La dynamique effective doit donc relier la courbure continue $G_{\mu\nu}[g^{\rm ent}]$
aux sources qui modifient la structure d'intrication. Ces sources ont deux origines.

\textit{Mati\`ere ordinaire.} Elle modifie localement l'\'etat quantique et donc
la connectivit\'e informationnelle du r\'eseau. Sa contribution appara\^{\i}t sous la forme
$\Geff(d_s)T_{\mu\nu}^{\rm matter}$, o\`u le couplage est modul\'e par la dimension
spectrale locale.

\textit{Gradients de dimension spectrale.} M\^eme en l'absence de source mat\'erielle
explicite, les gradients et variations de $d_s(x)$ contribuent \`a la courbure effective.
\`A l'ordre local minimal, la source g\'eom\'etrique construite \`a partir de $d_s$ doit
contenir les invariants covariants de plus bas ordre en nombre de d\'eriv\'ees :
$\nabla_\mu d_s\nabla_\nu d_s$ et $\nabla_\mu\nabla_\nu d_s$.
Ces termes d\'efinissent le tenseur effectif d'intrication $T_{\mu\nu}^{\ent}[d_s]$.

La forme effective minimale compatible avec la covariance et la localit\'e \`a grande
\'echelle est donc
\[
G_{\mu\nu}[g^{\rm ent}]
=
8\pi\Geff(d_s)\,T_{\mu\nu}^{\rm matter}
+
8\pi G\,T_{\mu\nu}^{\ent}[d_s].
\]
Cette \'equation n'est pas postul\'ee comme une loi microscopique exacte du r\'eseau.
Elle est la \textbf{fermeture continue minimale} de la cha\^{\i}ne
\[
\delta I\to\delta\ell\to\delta g_{\mu\nu}^{\rm ent}\to\delta G_{\mu\nu},
\qquad
\delta I\to\delta L\to\delta d_s\to\delta\Geff.
\]
Elle exprime que la relativit\'e g\'en\'erale \'emerge comme la r\'eponse
thermodynamique et continue du r\'eseau d'intrication aux variations locales
d'information mutuelle.

% ─────────────────────────────────────────────────────────────
\section{\'Equation d'Einstein effective \BuP}
\label{sec:einstein}
% ─────────────────────────────────────────────────────────────

Les sections pr\'ec\'edentes ont construit les deux objets g\'eom\'etriques
n\'ecessaires : la m\'etrique \'emergente $g_{\mu\nu}^{\rm ent}$, issue de la distance
d'intrication, et la dimension spectrale locale $d_s(x)$, issue du noyau de chaleur
du Laplacien d'intrication. Nous pr\'esentons maintenant l'\'equation de champ
effective et ses principales limites physiques.

Dans toute cette section, $\nabla_\mu$ d\'esigne la d\'eriv\'ee covariante associ\'ee
\`a la m\'etrique \'emergente $g_{\mu\nu}^{\rm ent}$.

\subsection{\'Equation centrale}

\[
\boxed{
G_{\mu\nu}\!\left[g^{\rm ent}\right]
=
8\pi\Geff(d_s)\,
T_{\mu\nu}^{\rm matter}
+
8\pi G\,
T_{\mu\nu}^{\ent}[d_s]
}
\]
o\`u
\[
\Geff(d_s)
=
\frac{2}{d_s-1}G.
\]

\subsection{Interpr\'etation des termes}

Le terme $8\pi\Geff(d_s)T_{\mu\nu}^{\rm matter}$ d\'ecrit la r\'eponse de la
g\'eom\'etrie \'emergente \`a la mati\`ere ordinaire, avec un couplage gravitationnel
modul\'e par la dimension spectrale locale.

Le terme $T_{\mu\nu}^{\ent}[d_s]$ ne repr\'esente pas une nouvelle composante
mat\'erielle. Il s'agit d'une \'ecriture effective des corrections g\'eom\'etriques
induites par les gradients et variations locales de la dimension spectrale.

\subsection{Tenseur effectif d'intrication}

\`A l'ordre effectif minimal compatible avec la covariance, la localit\'e effective et
l'invariance par renversement des coordonn\'ees spatiales, on peut \'ecrire
\[
T_{\mu\nu}^{\ent}[d_s]
=
\alpha
\!\left(
\nabla_\mu d_s\nabla_\nu d_s
-
\frac12 g_{\mu\nu}^{\rm ent}(\nabla d_s)^2
\right)
+
\beta
\!\left(
\nabla_\mu\nabla_\nu d_s
-
g_{\mu\nu}^{\rm ent}\Box d_s
\right).
\]

Le premier terme mesure l'\'energie effective associ\'ee aux gradients de dimension
spectrale. Il devient important lorsque $d_s(x)$ varie spatialement, par exemple dans
les r\'egions surdenses ou dans les zones de forte variation informationnelle.

Le second terme encode la r\'eponse de la g\'eom\'etrie aux variations secondes de
la dimension spectrale. Il joue le r\^ole d'un terme de rigidit\'e informationnelle :
il p\'enalise les variations rapides de $d_s$ et contr\^ole la fa\c{c}on dont la dimension
spectrale se lisse \`a grande \'echelle.

Les constantes $\alpha$ et $\beta$ ne sont pas introduites comme de nouveaux
param\`etres cosmologiques libres dans ce travail. Elles doivent \^etre d\'etermin\'ees,
en principe, par les fluctuations spectrales du Laplacien d'intrication dans les
simulations \BuP, ou contraintes ph\'enom\'enologiquement dans des applications
ult\'erieures.

\subsection{Conservation covariante et identit\'e de Bianchi}

La forme pr\'ec\'edente de $T_{\mu\nu}^{\ent}$ n'est cependant pas arbitraire :
elle doit \^etre compatible avec la conservation covariante impos\'ee par l'identit\'e
de Bianchi. Cette contrainte fixe la mani\`ere dont les variations de $d_s$ \'echangent
effectivement de l'\'energie-impulsion avec le secteur mati\`ere.

L'identit\'e de Bianchi impose $\nabla^\mu G_{\mu\nu}[g^{\rm ent}]=0$.
Par cons\'equent, le membre de droite de l'\'equation effective doit \^etre conserv\'e :
\[
\nabla^\mu
\left[
\Geff(d_s)\,T_{\mu\nu}^{\rm matter}
+
G\,T_{\mu\nu}^{\ent}[d_s]
\right]
=
0.
\]

Lorsque $d_s$ varie dans l'espace-temps, le couplage $\Geff(d_s)$ varie \'egalement,
et il peut exister un \'echange effectif entre le secteur mati\`ere et le secteur
g\'eom\'etrique d'intrication. On peut \'ecrire
\[
\nabla^\mu T_{\mu\nu}^{\rm matter}=Q_\nu,
\]
\[
\nabla^\mu T_{\mu\nu}^{\ent}
=
-\frac{1}{G}\,T_{\mu\nu}^{\rm matter}\,\nabla^\mu\Geff
-
\frac{\Geff}{G}\,Q_\nu,
\]
o\`u $Q_\nu$ d\'ecrit un \'echange effectif de quantit\'e d'\'energie-impulsion entre la
mati\`ere et la g\'eom\'etrie informationnelle.

Dans la limite o\`u $d_s$ est constant, on a $\nabla_\mu\Geff(d_s)=0$. Si l'on impose
alors $Q_\nu=0$, on retrouve la conservation usuelle :
$\nabla^\mu T_{\mu\nu}^{\rm matter}=0$ et $\nabla^\mu T_{\mu\nu}^{\ent}=0$.

\subsection{Limite relativit\'e g\'en\'erale}

La relativit\'e g\'en\'erale standard doit \^etre retrouv\'ee lorsque la g\'eom\'etrie
d'intrication devient homog\`ene et faiblement variable. Cette limite est d\'efinie par
\[
d_s(x)\rightarrow 3,
\qquad
\nabla_\mu d_s\rightarrow0,
\qquad
\nabla_\mu\nabla_\nu d_s\rightarrow0.
\]
Dans cette limite, $\Geff(d_s)\rightarrow G$ et $T_{\mu\nu}^{\ent}\rightarrow 0$.
L'\'equation de champ devient alors
\[
G_{\mu\nu}[g^{\rm ent}]
=
8\pi G\,T_{\mu\nu}^{\rm matter}.
\]
On retrouve la forme standard des \'equations d'Einstein, sans constante cosmologique
fondamentale ni composante d'\'energie sombre ajout\'ee.

\BuP\ ne remplace pas la relativit\'e g\'en\'erale dans les r\'egimes o\`u la dimension
spectrale est constante et proche de trois. La relativit\'e g\'en\'erale appara\^{\i}t
comme la limite homog\`ene, continue et thermodynamique d'une g\'eom\'etrie d'intrication
plus fondamentale.

\subsection{Limite cosmologique homog\`ene}

Dans la limite homog\`ene, l'\'equation effective \BuP\ contient comme cas particulier
l'\'equation de Friedmann modifi\'ee utilis\'ee dans les travaux cosmologiques
pr\'ec\'edents. En effet, lorsque
\[
d_s(x)\rightarrow d_{\rm bg}(z),
\]
la m\'etrique \'emergente prend la forme FLRW et le couplage gravitationnel devient
\[
\Geff(z)=\frac{2}{d_{\rm bg}(z)-1}G.
\]
Les composantes homog\`enes de l'\'equation effective donnent alors
\[
H^2(z)
=
\frac{8\pi G}{d_{\rm bg}(z)\left[d_{\rm bg}(z)-1\right]}
\rho_{\rm total}(z)
+
\Delta_{\rm ent}(z),
\]
o\`u $\Delta_{\rm ent}(z)$ regroupe les contributions temporelles issues de
$T_{\mu\nu}^{\ent}[d_s]$, c'est-\`a-dire les termes contenant $\dot d_{\rm bg}$
et $\ddot d_{\rm bg}$.

Dans la limite quasi-stationnaire, o\`u l'\'evolution de $d_{\rm bg}(z)$ est lente
sur l'\'echelle de Hubble, ces contributions peuvent \^etre n\'eglig\'ees ou absorb\'ees
dans la d\'efinition effective du fond cosmologique. On retrouve alors
\[
H^2(z)
=
\frac{8\pi G}{d_{\rm bg}(z)\left[d_{\rm bg}(z)-1\right]}
\rho_{\rm total}(z),
\]
qui est pr\'ecis\'ement l'\'equation de Friedmann modifi\'ee utilis\'ee dans le
mod\`ele cosmologique homog\`ene \BuP~\cite{hamdad2026bup}.

Cette \'equation a \'et\'e confront\'ee aux donn\'ees BAO et supernovæ dans les travaux
pr\'ec\'edents du programme \BuP. Le meilleur ajustement donne une dimension actuelle
$d_{\rm bg}(0)\simeq 2.40$, avec une am\'elioration statistique par rapport \`a $\LCDM$,
$\Delta{\rm BIC}=-3.71$.
L'acc\'el\'eration tardive n'est donc pas introduite par une constante cosmologique ou par
un fluide d'\'energie sombre : elle d\'ecoule de la dynamique g\'eom\'etrique de la
dimension effective lorsque $d_{\rm bg}(z)<3$ \`a bas redshift.

Le passage de $d_{\rm bg}>3$ \`a haut redshift vers $d_{\rm bg}<3$ \`a bas redshift
modifie simultan\'ement l'expansion du fond et la force gravitationnelle effective.
Ainsi, la limite homog\`ene de l'\'equation effective reproduit le r\'esultat principal
du mod\`ele cosmologique \BuP, tandis que les limites perturbative et locale reproduisent
respectivement les r\'esultats sur $\sig$ et l'effondrement non lin\'eaire des halos rares.

\subsection{Interpr\'etation sans secteur sombre}

Dans cette formulation, aucun secteur sombre fondamental n'est introduit. Il n'y a ni
constante cosmologique ind\'ependante, ni fluide d'\'energie sombre ajout\'e \`a la main.
L'acc\'el\'eration tardive est interpr\'et\'ee comme une cons\'equence g\'eom\'etrique
de l'\'evolution de la dimension spatiale effective $d_{\rm bg}(z)$.

Dans les travaux pr\'ec\'edents, l'expression \og\'energie noire\fg\ \'etait utilis\'ee
comme raccourci ph\'enom\'enologique pour d\'esigner l'acc\'el\'eration cosmique effective.
Dans le pr\'esent article, nous adoptons l'interpr\'etation g\'eom\'etrique stricte :
aucun fluide d'\'energie sombre fondamental n'est introduit.

La m\^eme fonction $d_{\rm bg}(z)$ agit simultan\'ement sur deux aspects de la
dynamique cosmologique. D'une part, elle modifie l'expansion du fond \`a travers
l'\'equation de Friedmann effective. D'autre part, elle modifie la force gravitationnelle
effective via $\Geff(z)=\frac{2}{d_{\rm bg}(z)-1}G$.

Ainsi, l'acc\'el\'eration tardive et la croissance modifi\'ee des structures ne sont pas
deux ph\'enom\`enes ind\'ependants : elles sont deux manifestations d'un seul m\'ecanisme
g\'eom\'etrique, \`a savoir la variation de la dimension effective de l'Univers.

\begin{table}[htbp]
\centering
\caption{Interpr\'etation g\'eom\'etrique des r\'egimes cosmologiques \BuP.}
\label{tab:regimes}
\smallskip
\begin{tabular}{lll}
\toprule
 & Haut redshift & Bas redshift \\
\midrule
Dimension effective $d_{\rm bg}$ & $\simeq 3.06>3$ & $\simeq 2.40<3$ \\
Couplage effectif $\Geff$ & $\simeq 0.97\,G$ & $\simeq 1.43\,G$ \\
Expansion & d\'ec\'el\'er\'ee & acc\'el\'er\'ee \\
Structures & croissance ralentie & collapse renforc\'e \\
Observable & $\sig$ r\'eduit & JWST / halos rares \\
\bottomrule
\end{tabular}
\end{table}

\subsection{Limite locale inhomog\`ene}

Dans une r\'egion surdense, la dimension spectrale locale re\c{c}oit une correction
informationnelle :
\[
d_s(x)
\simeq
d_{\rm bg}(z)-\varepsilon\delta(x).
\]
Le couplage gravitationnel local devient
\[
\Geff^{\rm local}
=
\frac{2}{d_{\rm bg}(z)-\varepsilon\delta-1}G.
\]
Au premier ordre en $\varepsilon\delta$,
\[
\frac{\Geff^{\rm local}}{\Geff^{\rm bg}}
\simeq
1+
\frac{\varepsilon\delta}{d_{\rm bg}(z)-1}.
\]
Ainsi, dans les r\'egions surdenses $(\delta>0)$, on obtient
$\Geff^{\rm local}>\Geff^{\rm bg}$.

Cette limite reproduit le m\'ecanisme d'effondrement non lin\'eaire local utilis\'e
pour d\'ecrire l'exc\`es de halos massifs \`a haut redshift~\cite{hamdad2026p4}.


% ─────────────────────────────────────────────────────────────
\section{Test numérique interne : corr\'elation courbure--dimension spectrale}
\label{sec:numerical-test}
% ─────────────────────────────────────────────────────────────

L'\'equation effective propos\'ee dans ce travail repose sur l'id\'ee que la
courbure \'emergente du graphe d'intrication est contr\^ol\'ee par les variations
locales de la dimension spectrale. Afin de tester cette relation au niveau discret,
nous avons construit un premier test num\'erique interne comparant un proxy de
courbure \'emergente \`a un proxy du tenseur spectral $T_{\mu\nu}^{\ent}[d_s]$.

Le point de d\'epart est une famille de matrices d'information mutuelle
$W_{ij}=I(i:j)$, g\'en\'er\'ees pour plusieurs tailles de r\'eseau
($N=9, 12, 16, 20$) et plusieurs valeurs du param\`etre de
non-localit\'e $\lambda$. Pour chaque matrice, nous construisons un graphe local
par $k$-plus proches voisins, les longueurs d'ar\^ete d'intrication
\[
\ell_{ij}
=
-\ell_0\log\!\left(
\frac{W_{ij}+\epsilon_I}{W_{\max}+\epsilon_I}
\right),
\]
puis la courbure discr\`ete d'Ollivier-Ricci $\kappa_{ij}$~\cite{ollivier2009}.

\`A partir de la courbure d'ar\^ete, nous d\'efinissons une courbure scalaire locale
moyenne $R_{\rm edge}$, puis plusieurs conventions discr\`etes pour un proxy du
tenseur d'Einstein \'emergent. Le scan num\'erique montre que la convention la
plus stable est
\[
G_{ij}^{\rm proxy}
=
-\kappa_{ij}
+
\frac{1}{2}\,R_{\rm edge}.
\]

En parall\`ele, la dimension spectrale locale $d_s(i)$ est extraite du noyau de
chaleur du Laplacien d'intrication. Nous construisons ensuite des proxys discrets
du tenseur spectral \`a partir des deux ingr\'edients
\[
(\nabla_{ij}d_s)^2
\qquad\text{et}\qquad
\Delta d_s,
\]
qui correspondent respectivement aux deux structures pr\'esentes dans
\[
T_{\mu\nu}^{\ent}[d_s]
=
\alpha
\!\left(
\nabla_\mu d_s\nabla_\nu d_s
-
\tfrac12 g_{\mu\nu}^{\rm ent}(\nabla d_s)^2
\right)
+
\beta
\!\left(
\nabla_\mu\nabla_\nu d_s
-
g_{\mu\nu}^{\rm ent}\Box d_s
\right).
\]

\subsection{R\'esultat sur $N=9, 12, 16$}

Sur l'ensemble des matrices $N=9, 12, 16$, le scan identifie une relation
dominante entre l'oppos\'e de la courbure scalaire d'ar\^ete et le laplacien
discret de la dimension spectrale~: $-R_{\rm edge}\sim\Delta d_s$.
La meilleure combinaison stable donne une corr\'elation de Pearson moyenne
$\langle r\rangle = 0.751$, avec une m\'ediane ${\rm med}(r)=0.892$, et une
corr\'elation de Spearman $\langle \rho\rangle=0.786$, avec un signe positif
dans $100\%$ des matrices test\'ees.

\subsection{R\'esultat sur $N=20$}

Pour les matrices $N=20$, la meilleure relation stable est
$G_{ij}^{\rm proxy}=-\kappa_{ij}+\frac{1}{2}R_{\rm edge}$,
corr\'el\'ee \`a un proxy spectral de la forme
$T_{ij}^{\rm proxy}=(\nabla_{ij}d_s)^2+|\Delta d_s|$,
donnant $\langle r\rangle=0.558$, ${\rm med}(r)=0.995$,
$\langle\rho\rangle=0.566$, avec un signe de Spearman positif dans
$100\%$ des matrices.

\begin{figure}[htbp]
  \centering
  \includegraphics[width=0.78\textwidth]{fig_best_G_vs_T_N20}
  \caption{Corr\'elation entre le proxy de courbure \'emergente
  $G_{ij}^{\rm proxy}=-\kappa_{ij}+\frac{1}{2}R_{\rm edge}$
  et le proxy spectral $T_{ij}^{\rm proxy}$ construit \`a partir de
  $(\nabla d_s)^2$ et $\Delta d_s$, pour les matrices $N=20$.}
  \label{fig:G-T-corr}
\end{figure}

\subsection{Robustesse en \'echelle locale $k$}

Nous avons test\'e la stabilit\'e du signal en variant l'\'echelle locale du
graphe avec $k=3, 5, 7$.

\begin{table}[htbp]
\centering
\caption{Robustesse du test $G_{ij}^{\rm proxy}$--$T_{ij}^{\rm proxy}$
pour $N=20$.}
\label{tab:einstein-corr-k}
\smallskip
\begin{tabular}{p{1cm}p{4.5cm}p{2cm}p{2cm}p{3cm}}
\toprule
$k$ & Proxy spectral dominant
    & $\langle r\rangle$
    & $\langle\rho\rangle$
    & Signe Spearman pos. \\
\midrule
3 & $(\nabla d_s)^2+|\Delta d_s|$ & 0.707 & 0.598 & 100\% \\
5 & $(\nabla d_s)^2+|\Delta d_s|$ & 0.558 & 0.566 & 100\% \\
7 & $(\nabla d_s)^2-\Delta d_s$   & 0.995 & 0.417 & 100\% \\
\bottomrule
\end{tabular}
\end{table}

\begin{figure}[htbp]
  \centering
  \includegraphics[width=0.90\textwidth]{fig_top10_stability_N20}
  \caption{Classement des combinaisons de proxys $G_{ij}^{\rm proxy}$ et
  $T_{ij}^{\rm proxy}$. Les meilleures combinaisons sont syst\'ematiquement
  construites \`a partir de la convention $-\kappa_{ij}+\frac{1}{2}R_{\rm edge}$
  et des variations locales de $d_s$.}
  \label{fig:proxy-ranking}
\end{figure}

La convention de courbure $G_{ij}^{\rm proxy}=-\kappa_{ij}+\frac{1}{2}R_{\rm edge}$
est stable pour toutes les valeurs de $k$ test\'ees. Les meilleurs proxys spectraux
sont syst\'ematiquement construits \`a partir de $(\nabla_{ij}d_s)^2$ et $\Delta d_s$.
Ce r\'esultat ne constitue pas encore une d\'erivation compl\`ete de l'\'equation
tensorielle continue, mais il fournit un premier test num\'erique interne de
coh\'erence : les proxys discrets de courbure \'emergente sont monotoniquement
corr\'el\'es aux proxys construits \`a partir des variations locales de $d_s$.
Cela soutient directement la structure propos\'ee pour $T_{\mu\nu}^{\ent}[d_s]$.

% ─────────────────────────────────────────────────────────────
\section{Discussion et perspectives}
\label{sec:discussion}
% ─────────────────────────────────────────────────────────────

L'\'equation effective propos\'ee dans ce travail fournit un cadre g\'eom\'etrique unique
permettant de relier les diff\'erents r\'esultats cosmologiques obtenus dans le
programme \BuP. Son int\'er\^et principal est de montrer que les corrections introduites
pr\'ec\'edemment ne sont pas des prescriptions ind\'ependantes, mais peuvent \^etre
interpr\'et\'ees comme des limites particulières d'une m\^eme dynamique \'emergente.

\subsection{Unification des r\'egimes \BuP}

Paper 5 n'ajoute pas un nouveau mod\`ele : il unifie les trois r\'egimes d\'ej\`a obtenus.

\begin{table}[htbp]
\centering
\caption{Les trois r\'egimes cosmologiques pr\'ec\'edents comme limites de l'\'equation
effective \BuP.}
\label{tab:unification}
\smallskip
\begin{tabular}{p{2.4cm}p{2.5cm}p{3.4cm}p{4.6cm}}
\toprule
Travail & R\'egime & Relation cl\'e & Interpr\'etation dans Paper 5 \\
\midrule
\cite{hamdad2026bup} &
Homog\`ene &
$d_s\rightarrow d_{\rm bg}(z)$ &
Limite FLRW de l'\'equation effective \\[4pt]
\cite{hamdad2026p3} &
Lin\'eaire &
$\Geff(z)=2G/[d_{\rm bg}(z)-1]$ &
Couplage gravitationnel modul\'e par $d_s$ \\[4pt]
\cite{hamdad2026p4} &
Non lin\'eaire local &
$d_s=d_{\rm bg}-\varepsilon\delta$ &
Correction inhomog\`ene de la dimension spectrale \\
\bottomrule
\end{tabular}
\end{table}

Dans la limite homog\`ene, la dimension spectrale devient $d_s(x)\rightarrow d_{\rm bg}(z)$.
L'\'equation effective se r\'eduit \`a une dynamique de fond dans laquelle $\Geff(z)$
d\'epend de la dimension \'emergente. Cette limite correspond au mod\`ele cosmologique
\`a dimension variable contraint par les donn\'ees BAO, supernovæ et CMB.

Dans le r\'egime lin\'eaire, la m\^eme d\'ependance de $\Geff$ \`a $d_s$ modifie la
croissance des perturbations. Lorsque $d_{\rm bg}(z)\rightarrow\DC>3$ \`a haut redshift,
on obtient $\Geff<G$, ce qui produit une l\'eg\`ere suppression de la croissance des
structures et permet de retrouver $\sig\simeq0.772$.

Enfin, dans les r\'egions localement surdenses, la dimension spectrale re\c{c}oit une
correction $d_s\simeq d_{\rm bg}(z)-\varepsilon\delta(x)$, augmentant le couplage
local $\Geff^{\rm local}>\Geff^{\rm bg}$ et acc\'el\'erant l'effondrement des pics
rares. Le r\'egime non lin\'eaire utilis\'e pour interpr\'eter l'exc\`es JWST appara\^{\i}t
donc comme une limite locale de la m\^eme \'equation effective.

La dimension de fond contr\^ole la moyenne cosmologique ; la dimension locale contr\^ole
les fluctuations rares. Cette s\'eparation explique pourquoi la r\'esolution de la
tension $\sig$ peut coexister avec un exc\`es de structures massives \`a haut redshift.

\subsection{Diff\'erence avec les th\'eories scalaire-tenseur}

La forme du tenseur d'intrication $T_{\mu\nu}^{\ent}[d_s]$ ressemble formellement
\`a celle d'un secteur scalaire effectif. Cependant, l'interpr\'etation physique est
diff\'erente. Dans une th\'eorie scalaire-tenseur classique, le champ scalaire est
introduit comme un degr\'e de libert\'e fondamental suppl\'ementaire. Dans \BuP,
$d_s(x)$ n'est pas ajout\'e \`a la th\'eorie : il est d\'efini \`a partir du spectre
du Laplacien d'intrication. Il s'agit donc d'une observable g\'eom\'etrique \'emergente,
non d'une nouvelle mati\`ere exotique.

\subsection{Limites actuelles}

La pr\'esente formulation reste effective. Les constantes $\alpha$ et $\beta$ doivent
encore \^etre d\'etermin\'ees depuis les simulations \BuP\ \`a grande taille. La limite
continue du graphe d'intrication reste une conjecture. La compatibilit\'e avec les
tests locaux de la relativit\'e g\'en\'erale (syst\`eme solaire, pulsars binaires,
ondes gravitationnelles) doit \^etre quantifi\'ee.

\subsection{Pr\'edictions observationnelles}

(i) \textit{Amas de galaxies} : $\Geff^{\rm local}>\Geff^{\rm bg}$ dans les r\'egions
denses pourrait affecter les comptages d'amas \`a $z\sim1$--$2$ (Euclid, eROSITA, LSST).

(ii) \textit{Lentillage gravitationnel} : les profils de lentillage pourraient recevoir
des corrections li\'ees aux gradients de $d_s$.

(iii) \textit{Formation pr\'ecoce des galaxies} : le m\'ecanisme $d_s=d_{\rm bg}-\varepsilon\delta$
pr\'evoit un effet pr\'ef\'erentiel sur la queue rare, testable avec Roman Space Telescope
et Euclid.

(iv) \textit{Ondes gravitationnelles} : $T_{\mu\nu}^{\ent}$ pourrait induire des
corrections \`a la propagation des perturbations m\'etriques cosmologiques.

\subsection{Programme num\'erique}

Le test d\'ecisif consistera \`a reconstruire num\'eriquement depuis un \'etat \BuP\ :
\[
W_{ij}\to\ell_{ij}\to d_{ij}^{\rm ent}\to g_{\mu\nu}^{\rm ent}\to R_{\mu\nu}^{\rm ent},
\]
puis \`a v\'erifier la relation effective
$G_{\mu\nu}[g^{\rm ent}]\sim 8\pi\Geff(d_s)T_{\mu\nu}^{\rm matter}+8\pi GT_{\mu\nu}^{\ent}$.

% ─────────────────────────────────────────────────────────────
\section{Conclusion}
\label{sec:conclusion}
% ─────────────────────────────────────────────────────────────

Nous avons propos\'e une \'equation d'Einstein effective \'emergente dans le cadre
\BuP. Le point de d\'epart est que la g\'eom\'etrie n'est pas postul\'ee comme
fondamentale : elle est reconstruite \`a partir de la structure d'intrication
d'un \'etat quantique global.

La cha\^{\i}ne g\'eom\'etrique centrale est
\[
I(i:j)
\rightarrow
\ell_{ij}
\rightarrow
d_{ij}^{\rm ent}
\rightarrow
g_{\mu\nu}^{\rm ent}
\rightarrow
d_s(x)
\rightarrow
\Geff(d_s).
\]
Elle relie directement l'information mutuelle microscopique \`a la m\'etrique \'emergente,
\`a la dimension spectrale locale et au couplage gravitationnel effectif.

L'\'equation propos\'ee est
\[
G_{\mu\nu}[g^{\rm ent}]
=
8\pi\Geff(d_s)\,T_{\mu\nu}^{\rm matter}
+
8\pi G\,T_{\mu\nu}^{\ent}[d_s],
\]
avec
\[
\Geff(d_s)=\frac{2}{d_s-1}G,
\quad
T_{\mu\nu}^{\ent}[d_s]
=\alpha\!\left(\nabla_\mu d_s\nabla_\nu d_s-\tfrac12 g_{\mu\nu}^{\rm ent}(\nabla d_s)^2\right)
+\beta\!\left(\nabla_\mu\nabla_\nu d_s-g_{\mu\nu}^{\rm ent}\Box d_s\right).
\]

Aucune constante cosmologique $\Lambda$ n'est postul\'ee. L'acc\'el\'eration tardive
est une cons\'equence g\'eom\'etrique de $d_{\rm bg}(z)<3$.

Cette \'equation poss\`ede trois limites physiques importantes.

Premi\`erement, lorsque $d_s(x)\rightarrow 3$, $\nabla_\mu d_s\rightarrow0$, elle
redonne la relativit\'e g\'en\'erale standard sans constante cosmologique fondamentale.

Deuxi\`emement, dans la limite homog\`ene $d_s(x)\rightarrow d_{\rm bg}(z)$, elle
reproduit la cosmologie \`a dimension variable et le couplage
$\Geff(z)=\frac{2}{d_{\rm bg}(z)-1}G$.

Troisi\`emement, dans les r\'egions surdenses $d_s(x)\simeq d_{\rm bg}(z)-\varepsilon\delta(x)$,
elle redonne le m\'ecanisme local d'amplification gravitationnelle utilis\'e pour
expliquer l'effondrement pr\'ecoce des halos rares.

Ainsi, les r\'esultats cosmologiques pr\'ec\'edents du programme \BuP\ apparaissent
comme des limites particulières d'une m\^eme structure g\'eom\'etrique. La relativit\'e
g\'en\'erale n'est pas abandonn\'ee : elle \'emerge comme la limite continue, homog\`ene
et thermodynamique d'un r\'eseau d'intrication quantique plus fondamental.

Dans cette formulation, les effets g\'en\'eralement attribu\'es \`a l'\'energie sombre
sont r\'einterpr\'et\'es comme des effets de g\'eom\'etrie \'emergente, et non comme
une nouvelle composante du contenu \'energ\'etique de l'Univers. Les prochaines \'etapes
consisteront \`a d\'eterminer les constantes effectives $\alpha$ et $\beta$ \`a partir
des simulations \BuP, \`a tester la limite newtonienne, et \`a confronter les
pr\'edictions aux donn\'ees de lentillage, d'amas de galaxies et de formation pr\'ecoce
des structures.

% ─────────────────────────────────────────────────────────────
\bibliographystyle{plain}
\begin{thebibliography}{99}

\bibitem{hamdad2026bup}
F.~Hamdad,
\textit{Contraintes cosmologiques sur une dimension spatiale \'emergente :
dimension variable, \'energie noire et $\sigma_8$ dans le cadre
Bottom-Up Quantum Gravity},
Pr\'epublication HAL hal-05590614v1 (2026),
\url{https://hal.science/hal-05590614v1}.

\bibitem{hamdad2026p3}
F.~Hamdad,
\textit{R\'esolution g\'eom\'etrique de la tension $\sigma_8$ par une
dimension cosmologique \'emergente},
Pr\'epublication (2026).

\bibitem{hamdad2026p4}
F.~Hamdad,
\textit{Effondrement non-lin\'eaire par r\'eduction dimensionnelle locale
dans la Gravit\'e Quantique Bottom-Up},
Pr\'epublication (2026).

\bibitem{jacobson1995}
T.~Jacobson,
\textit{Thermodynamics of spacetime: the Einstein equation of state},
Phys.\ Rev.\ Lett.\ \textbf{75}, 1260 (1995).

\bibitem{ollivier2009}
Y.~Ollivier,
\textit{Ricci curvature of Markov chains on metric spaces},
J.\ Funct.\ Anal.\ \textbf{256}, 810--864 (2009).

\bibitem{ambjorn2005}
J.~Ambjorn, J.~Jurkiewicz, R.~Loll,
\textit{Spectral dimension of the universe},
Phys.\ Rev.\ Lett.\ \textbf{95}, 171301 (2005).

\bibitem{penrose2003}
M.~Penrose,
\textit{Random Geometric Graphs},
Oxford University Press (2003).

\bibitem{planck2018}
Planck Collaboration,
\textit{Planck 2018 results VI. Cosmological parameters},
Astron.\ Astrophys.\ \textbf{641}, A6 (2020).

\bibitem{desi2024}
DESI Collaboration,
\textit{DESI 2024 VI: Cosmological Constraints from Baryon Acoustic
Oscillations},
arXiv:2404.03002 (2024).

\bibitem{labbe2023}
I.~Labb\'e \textit{et al.},
\textit{A population of red candidate massive galaxies $\sim$600~Myr
after the Big Bang},
Nature \textbf{616}, 266 (2023).

\end{thebibliography}

\end{document}
